\section*{Chapitre \ref{ch:b1:heat-engines} --- Moteurs et machines thermiques}

\begin{solution}{exo:b1:heat-engines:1}
$\eta_C = 1 - 300/800 = 0.625$. À $0.40$~: $\dot Q_h = 1.0/0.40 = \qty{2.5}{GW}$
prélevés au combustible, $\dot Q_c = 2.5 - 1.0 = \qty{1.5}{GW}$ vers la rivière ---
on jette plus de chaleur qu'on n'en vend.
\end{solution}

\begin{solution}{exo:b1:heat-engines:2}
$e_f \leq 275/(298 - 275) = 12$. Avec $e_f = 3$~: $P = 100/3 = \qty{33}{W}$
(et \qty{133}{W} partent dans la cuisine).
\end{solution}

\begin{solution}{exo:b1:heat-engines:3}
$r = 9$~: $9^{0.4} = 2.41$, $\eta = 0.585$~; $r = 11$~: $11^{0.4} = 2.61$,
$\eta = 0.62$. À $r = 20$, le mélange essence--air, comprimé jusqu'à
\qty{900}{K} environ, s'enflammerait avant l'étincelle (cliquetis) et
détruirait le moteur~; seul le Diesel, qui ne comprime que de l'air, peut y aller.
\end{solution}

\begin{solution}{exo:b1:heat-engines:4}
$e_p \leq 300/20 = 15$~: puissance idéale $5/15 = \qty{0.33}{kW}$~; réelle, $5/3.5 =
\qty{1.43}{kW}$~; résistance, \qty{5}{kW}.
\end{solution}

\begin{solution}{exo:b1:heat-engines:5}
Adiabatiques~: $V_C = V_B(T_h/T_c)^{1/(\gamma-1)} = 3.0 \times (5/3)^{2.5} =
\qty{10.8}{L}$, $V_D = 1.0 \times 3.59 = \qty{3.59}{L}$ (donc $V_C/V_D = 3 =
V_B/V_A$). $Q_h = RT_h\ln3 = 8.314 \times 500 \times 1.099 = \qty{4.57}{kJ}$,
$Q_c = -RT_c\ln3 = \qty{-2.74}{kJ}$, $W = -(4.57 - 2.74) = \qty{-1.83}{kJ}$,
$\eta = 1.83/4.57 = 0.40 = 1 - 300/500$.
\end{solution}

\begin{solution}{exo:b1:heat-engines:6}
Première annonce~: l'énergie s'équilibre, $200 = 1000 - 800$~; Clausius $1000/400 - 800/300
= 2.5 - 2.67 = \qty{-0.17}{J/K}$, donc $S_{\text{créée}} = \qty{0.17}{J/K} \geq 0$~:
c'est possible ($\eta = 0.20$). Seconde~: $1000/400 - 700/300 = +\qty{0.17}{J/K} > 0$~:
il faudrait détruire de l'entropie --- impossible. Maximum~: $\eta_C = 1 -
300/400 = 0.25$, soit \qty{250}{J}.
\end{solution}

\begin{solution}{exo:b1:heat-engines:7}
$r^{\gamma-1} = 18^{0.4} = 3.18$~; $\rho^\gamma = 2^{1.4} = 2.64$~: $\eta = 1 -
1.64/(1.4 \times 3.18 \times 1) = 0.63$. Otto à $r = 18$~: $1 - 1/3.18 = 0.69$~;
à $r = 9$~: $0.585$. Le Diesel ne comprime que de l'air~: aucun carburant n'est
là pour s'enflammer trop tôt, et celui que l'on injecte dans un air à \qty{900}{K}
brûle à mesure qu'il entre --- la limite du cliquetis n'existe pas.
\end{solution}

\begin{solution}{exo:b1:heat-engines:8}
$Q_h = RT_h\ln2 = 8.314 \times 600 \times 0.693 = \qty{3.46}{kJ}$, $Q_c = -8.314
\times 300 \times 0.693 = \qty{-1.73}{kJ}$~: $W = \qty{-1.73}{kJ}$ par cycle. Avec
le régénérateur, $\eta = 1.73/3.46 = 0.50 = 1 - 300/600$. Sans lui, la source
chaude doit aussi fournir $C_{V,m}(T_h - T_c) = 2.5 \times 8.314 \times 300 =
\qty{6.24}{kJ}$~: $\eta = 1.73/(3.46 + 6.24) = 0.18$.
\end{solution}

\begin{solution}{exo:b1:heat-engines:9}
Entropie des deux blocs~: $\Delta S = C\ln(T_f/T_1) + C\ln(T_f/T_2) =
C\ln\bigl(T_f^2/T_1T_2\bigr) = S_{\text{créée}} \geq 0$ (l'entropie propre du
moteur revient à sa valeur à chaque cycle), donc $T_f \geq \sqrt{T_1T_2}$, et
$W = C(T_1 + T_2 - 2T_f)$ (premier principe) est maximal pour le plus petit $T_f$~:
le cas réversible, $T_f = \sqrt{400 \times 300} = \qty{346.4}{K}$,
$W_{\max} = 4000 \times (700 - 692.8) = \qty{29}{kJ}$. Contact direct~: $T_f =
\qty{350}{K}$, aucun travail, $S_{\text{créée}} = 4000\ln(350^2/120000) =
\qty{82}{J/K}$.
\end{solution}

\begin{solution}{exo:b1:heat-engines:10}
(a) $P = \dot Q_h - |\dot Q_c| = K(T_h - x) - K(y - T_c)$~; le cœur réversible
n'échange aucune entropie nette~: $K(T_h - x)/x = K(y - T_c)/y$. (b) En croisant~:
$y(T_h - x) = x(y - T_c)$, donc $y(T_h - 2x) = -xT_c$ et $y = xT_c/(2x - T_h)$.
Puis $P = K(T_h - x)(1 - y/x) = K(T_h - x)(2x - T_h - T_c)/(2x - T_h)$~; avec
$u = 2x - T_h$, $T_h - x = (T_h - u)/2$ et $P = K(T_h - u)(u - T_c)/(2u) =
\tfrac{K}{2}[(T_h + T_c) - u - T_hT_c/u]$. (c) $\dd P/\dd u = \tfrac K2(-1 +
T_hT_c/u^2) = 0$ en $u = \sqrt{T_hT_c}$, qui est un maximum~; alors $x = (T_h +
\sqrt{T_hT_c})/2$, $y = xT_c/u = (\sqrt{T_hT_c} + T_c)/2$, et $\eta^\ast = 1 - y/x
= 1 - \sqrt{T_c}(\sqrt{T_h} + \sqrt{T_c})/[\sqrt{T_h}(\sqrt{T_h} + \sqrt{T_c})] = 1 -
\sqrt{T_c/T_h}$~; $P_{\max} = \tfrac K2(\sqrt{T_h} - \sqrt{T_c})^2$. (d)
$\eta^\ast = 1 - \sqrt{0.375} = 0.39$, contre $\eta_C = 0.625$~: le $0.40$ réel
de l'\cref{exo:b1:heat-engines:1} est un moteur construit pour la puissance,
non pour le \omterm{def:b1:heat-engines:efficiency}{rendement}.
\end{solution}

\begin{solution}{exo:b1:heat-engines:11}
Cycle~: $Q_h + Q_c + Q_0 = 0$ et $Q_h/T_h + Q_c/T_c + Q_0/T_0 \leq 0$.
Éliminons $Q_0 = -(Q_h + Q_c)$~: $Q_h(1/T_h - 1/T_0) + Q_c(1/T_c - 1/T_0) \leq 0$,
c'est-à-dire $Q_c(T_0 - T_c)/(T_cT_0) \leq Q_h(T_h - T_0)/(T_hT_0)$, d'où
$Q_c/Q_h \leq (1 - T_0/T_h)\,T_c/(T_0 - T_c)$. Le premier facteur est le
\omterm{def:b1:heat-engines:efficiency}{rendement} d'un moteur de Carnot entre le brûleur et la pièce, le second
l'\omterm{def:b1:heat-engines:efficiency}{efficacité} d'un \omterm{def:b1:heat-engines:machine}{réfrigérateur} de Carnot entre la chambre et la pièce~: la
borne est ce que donne le couplage idéal des deux. Numériquement~: $0.25 \times 11 = 2.75$
(les \omterm{def:b1:heat-engines:machine}{réfrigérateurs} à absorption réels~: environ $0.6$).
\end{solution}

\begin{solution}{exo:b1:heat-engines:12}
(a) \qty{0.40}{J}. (b) $0.40 \times 3.5 = \qty{1.4}{J}$. (c) \qty{0.90}{J}.
(d) Moteur $\eta = 1 - 293/800 = 0.634$, pompe $e_p = 293/20 = 14.65$~:
$0.634 \times 14.65 = \qty{9.3}{J}$ (et les \qty{0.37}{J} que le moteur
rejette à \qty{293}{K} pourraient aussi chauffer la maison~: \qty{9.7}{J}). Plus
d'un joule, parce que le reste est prélevé à l'air extérieur~: l'énergie est
conservée, et le second principe demande seulement que l'entropie ne soit pas
détruite, ce qu'une chaîne réversible respecte exactement.
\end{solution}

\begin{solution}{pb:b1:heat-engines:1}
\textbf{1.} $P = 250 \times 25 = \qty{6.25}{kW}$.

\textbf{2.} Perte moyenne $250 \times 13 = \qty{3.25}{kW}$ pendant $200 \times 86400 =
\qty{1.73e7}{s}$~: $E = \qty{5.6e10}{J} = \qty{15600}{kW h}$.

\textbf{3.} $15600 \times 0.25 = 3900$~\texteuro.

\textbf{4.} Gaz $15600/0.90 = \qty{17300}{kW h}$~: $1730$~\texteuro.

\textbf{5.} Résistances~: $15600/0.40 = \qty{39000}{kW h}$ de combustible~;
chaudière \qty{17300}{kW h}~: $2.25$ fois moins. L'électricité est du travail,
la forme qui peut tout faire~; la transformer en chaleur à
\qty{20}{\degreeCelsius} dans une résistance jette à la fois les $60\%$ du
combustible déjà perdus à la centrale et la possibilité de pomper de la chaleur.

\textbf{6.} Sur un cycle, $W + Q_h + Q_c = 0$ et $Q_h/T_{\mathrm{int}} +
Q_c/T_{\mathrm{ext}} \leq 0$ avec $Q_h < 0$, $Q_c > 0$~: $Q_c \leq
|Q_h|T_{\mathrm{ext}}/T_{\mathrm{int}}$, donc $W = |Q_h| - Q_c \geq |Q_h|(1 -
T_{\mathrm{ext}}/T_{\mathrm{int}})$ et $e_p = |Q_h|/W \leq T_{\mathrm{int}}/(T_{\mathrm{int}} -
T_{\mathrm{ext}})$.

\textbf{7.} $293/25 = 11.7$~; $6.25/11.7 = \qty{0.53}{kW}$.

\textbf{8.} $293/13 = 22.5$~; $15600/22.5 = \qty{690}{kW h}$ (environ
$170$~\texteuro).

\textbf{9.} Le travail par joule remonté vaut $(T_{\mathrm{int}} - T_{\mathrm{ext}})/
T_{\mathrm{int}}$, la «~hauteur~» de la remontée~; plus il fait froid dehors,
plus haut chaque joule doit être pompé --- et plus il en faut, puisque la perte
est elle aussi proportionnelle à $T_{\mathrm{int}} - T_{\mathrm{ext}}$~: la
\omterm{thm:b1:dc-circuits:kirchhoff}{puissance électrique} croît comme le carré de l'écart de température, au pire
précisément quand le chauffage compte le plus.

\textbf{10.} De l'air extérieur, que l'on refroidit~: $Q_c = 6.25 - 0.53 =
\qty{5.7}{kW}$. Le second principe interdit à la chaleur d'aller
\emph{spontanément} du froid vers le chaud~; ici, le travail paie la remontée.
Vérification~: entropie prise à l'air $5.72/268 = \qty{21.3}{W/K}$, entropie
cédée à la maison $6.25/293 =
\qty{21.3}{W/K}$ --- rien de détruit, rien de créé~: la pompe idéale.

\textbf{11.} Compresseur (reçoit $W$), condenseur (cède $|Q_h|$), détendeur
(ni l'un ni l'autre), évaporateur (reçoit $Q_c$). La chaleur ne descend que la
pente des températures~: le fluide doit être plus froid que l'air pour lui
prendre de la chaleur, et plus chaud que l'eau de la maison pour lui en céder.

\textbf{12.} $T_{\mathrm{ev}} = 268 - 8 = \qty{260}{K}$~; $e_{p,C} = 313/(313 - 260)
= 5.9$.

\textbf{13.} $e_p = 0.55 \times 5.9 = 3.25$~; $W = 6.25/3.25 = \qty{1.92}{kW}$~;
$Q_c = 6.25 - 1.92 = \qty{4.33}{kW}$ prélevés à l'air.

\textbf{14.} $\dot m = 4330/200000 = \qty{22}{g/s}$, soit \qty{1.3}{kg} par minute.

\textbf{15.} $T_{\mathrm{ev}} = \qty{272}{K}$, $313/41 = 7.6$, $e_p = 4.2$~;
$W = 3.25/4.2 = \qty{0.77}{kW}$.

\textbf{16.} $e_p = 0.55 \times 313/(321 - T_{\mathrm{ext}}) = 172/(321 -
T_{\mathrm{ext}})$ ($T_{\mathrm{ext}}$ en \omterm{def:b1:kinetic-theory:temperature}{kelvins}). Sur la saison~: $15600/4.2 =
\qty{3700}{kW h}$, $930$~\texteuro\ --- contre $3900$ avec les résistances et
$1730$ avec le gaz.

\textbf{17.} $3700/0.40 = \qty{9300}{kW h}$ de combustible~: $46\%$ de moins que
la chaudière et quatre fois moins que les résistances, et cela au travers d'une
centrale à $40\%$.

\textbf{18.} Besoin $250 \times 35 = \qty{8.75}{kW}$~; $T_{\mathrm{ev}} = \qty{250}{K}$,
$313/63 = 5.0$, $e_p = 2.7$, $W = \qty{3.2}{kW}$ si elle le pouvait. Capacité
\qty{3.8}{kW}~: appoint $8.75 - 3.8 \approx \qty{5}{kW}$ de résistances. Une pompe
dimensionnée pour l'heure la plus froide serait surdimensionnée --- et
fonctionnerait par cycles courts, sans \omterm{def:b1:heat-engines:efficiency}{rendement} --- tout l'hiver~; on la
dimensionne à la température de base et l'on laisse une résistance bon marché
couvrir l'extrême rare.

\textbf{19.} $338/(338 - 260) = 4.33$, $e_p = 2.4$, $W = 6.25/2.4 =
\qty{2.6}{kW}$~: $36\%$ d'électricité de plus qu'avec le plancher à
\qty{35}{\degreeCelsius}. Les émetteurs à basse température sont la moitié de
l'installation.

\textbf{20.} Entropie sur un cycle~: $Q_h/T_{\mathrm{int}} + Q_c/T_{\mathrm{ext}} +
S_{\text{créée}} = 0$ donne $Q_c = T_{\mathrm{ext}}|Q_h|/T_{\mathrm{int}} -
T_{\mathrm{ext}}S_{\text{créée}}$~; on reporte dans $W = |Q_h| - Q_c$.

\textbf{21.} Condenseur~: $6250\,(1/293 - 1/313) = \qty{1.36}{W/K}$~; évaporateur~:
$4330\,(1/260 - 1/268) = \qty{0.50}{W/K}$.

\textbf{22.} $268 \times 1.86 = \qty{0.50}{kW}$~: le quart des
\qty{1.92}{kW}.

\textbf{23.} $(1.92 - 0.53)/268 = \qty{5.2}{W/K}$~; les échangeurs en font
$36\%$~; le reste vient du compresseur (frottement, compression non
isentropique, pertes du moteur), du laminage dans le détendeur (une détente
libre), des pertes de charge et des ventilateurs.

\textbf{24.} $T_{\mathrm{cd}} = \qty{303}{K}$, $T_{\mathrm{ev}} = \qty{264}{K}$~:
condenseur $6250\,(1/293 - 1/303) = \qty{0.70}{W/K}$, évaporateur
$4330\,(1/264 - 1/268) = \qty{0.25}{W/K}$~: \qty{0.95}{W/K} au lieu de
$1.86$, soit $268 \times 0.91 \approx \qty{0.25}{kW}$ économisés ($13\%$~; la
machine entière, calée à $0.55$ d'une meilleure valeur de Carnot, gagnerait
davantage). Diviser un écart par deux, c'est doubler la surface d'échange~:
des échangeurs plus gros, plus chers, et plus de puissance de ventilation.

\textbf{25.} $15600/3700 = 4.2$ kilowattheures de chaleur par kilowattheure
d'électricité~; $930$~\texteuro\ contre $1730$ pour le gaz~; carbone~: $3700
\times 0.060 = \qty{220}{kg}$ contre $17300 \times 0.20 = \qty{3500}{kg}$ ---
quinze fois moins.
\end{solution}
